\chapter{Megvalósítás}
\label{ch:megvalositas}
% Terjedelem: kb. 9–10 oldal (~15 000 karakter)
% Kódrészletből csak a lényeges, rövid részeket! A teljes kód mellékletbe / repóba.

\section{A backend felépítése}
% Mappaszerkezet: api/ (auth, exams, attempts, results), core/ (database, security, timeutils),
% models/, schemas/ (Pydantic), services/ (grading, sandbox, socket_server).
\todo{Szöveg}

\section{Hitelesítés és jogosultságkezelés}
% bcrypt hash, JWT (HS256, 1 nap), oktatói meghívókód, szerepkör-ellenőrzés.
\todo{Szöveg}

\section{A vizsgafolyamat és az időkezelés}
% Kezdés -> határidő = min(kezdés + időtartam, vizsga záró időpontja); 30 mp-es türelmi idő;
% automatikus mentés; lejáratkor automatikus lezárás (finalize_if_expired).
\todo{Szöveg}

\section{Az automatikus pontozás}
% Típusonként. A több jó válaszos feladat képlete:
A több helyes választ is tartalmazó feladatoknál a rendszer részpontot ad, a hibás jelöléseket
pedig levonja:
\begin{equation}
  P = \max\!\left(0,\ \frac{|S \cap H| - |S \setminus H|}{|H|}\right) \cdot P_{\max},
  \label{eq:tobbvalasz}
\end{equation}
ahol $S$ a megjelölt, $H$ a helyes válaszok halmaza, $P_{\max}$ pedig a feladat pontértéke.
% Rövid válasz: kis/nagybetű- és szóközfüggetlen, számoknál 3,14 = 3.14.
% Kódfeladat: pont = sikeres tesztek / összes teszt * pont.
\todo{Szöveg}

\section{Biztonságos kódfuttatás Docker-sandboxban}
% Egy konténer / kiértékelés, runner script, nobody felhasználó, kimenet- és időkorlát.
A konténerre alkalmazott korlátozásokat a \ref{tab:sandbox}.~táblázat, a konténer indítását
a \ref{lst:sandbox}.~kódrészlet mutatja be.

\begin{table}[H]
  \centering
  \caption{A sandbox konténer korlátozásai}
  \label{tab:sandbox}
  \begin{tabularx}{\textwidth}{@{}lX@{}}
    \toprule
    Korlátozás & Érték, cél \\
    \midrule
    Hálózat             & nincs (\texttt{network\_mode="none"}) – adatszivárgás ellen \\
    Memória             & 128 MB \\
    CPU                 & 0,5 mag \\
    Folyamatok száma    & legfeljebb 64 (fork bomb ellen) \\
    Fájlrendszer        & csak olvasható; \texttt{/tmp}: 16 MB-os tmpfs \\
    Jogosultságok       & minden capability eldobva, \texttt{no-new-privileges}; a kód \texttt{nobody}-ként fut \\
    Időkorlát           & 5 mp tesztesetenként \\
    Kimenet             & legfeljebb 64 KB tesztesetenként \\
    \bottomrule
  \end{tabularx}
\end{table}

\begin{lstlisting}[language=Python,caption={A sandbox konténer indítási paraméterei (\texttt{sandbox.py}, részlet)},label={lst:sandbox}]
run_kwargs = dict(
    image=SANDBOX_IMAGE,
    command=["python", "-c", _RUNNER],
    network_mode="none",
    mem_limit="128m",
    nano_cpus=500_000_000,
    pids_limit=64,
    read_only=True,
    tmpfs={"/tmp": "size=16m,mode=1777"},
    cap_drop=["ALL"],
    cap_add=["SETUID", "SETGID", "KILL"],
    security_opt=["no-new-privileges"],
)
\end{lstlisting}

\section{Valós idejű felügyelet}
% Kliens: useProctoring hook (blur + visibilitychange, dupla számolás kiszűrése).
% Szerver: Socket.IO események (connect JWT-vel, join_exam, report_violation, exam_finished),
% szobák vizsgánként, a szabálysértések az adatbázisban tárolódnak.
\todo{Szöveg}

\section{A felhasználói felület}
% Képernyőképek (abrak/ mappába): bejelentkezés, hallgatói műszerfal, vizsgaoldal kódszerkesztővel,
% oktatói feladatszerkesztő, élő monitor, eredmények. Minden képre hivatkozz a szövegben!
\todo{Szöveg és képernyőképek}
